\section{Introduction}\label{sec:intro}

Okun's law is the empirical regularity that unemployment falls when output grows faster
than its trend and rises when it grows more slowly \citep{okun1962}. It is among the most
cited relations in macroeconomics and one of the least settled: its coefficient varies
across countries and periods, and its meaning depends on what moves output in the first
place \citep{ball2017okun}.

\paragraph{Motivation.} This study was prompted by a reported hierarchical Bayesian analysis
of the \eu\ for 2009--2026. That analysis described Spain as having a large association
between GDP growth and unemployment, the Netherlands as having a delayed response, and
Luxembourg as having a weak one, and it suggested measurement and institutional
explanations. We test two explanations of that kind: a mismatch between production by
cross-border commuters and resident unemployment in Luxembourg, and employment
protection. We do not have that analysis's code, data vintage or numerical
estimates. We therefore treat its statements as hypotheses and examine them
independently. The three descriptive claims were written into the protocol of the
Bayesian stage before its estimates, and the two explanations into the protocol of the
case study (Appendix~\ref{app:prereg}). Portugal, about which no claim was made, is
included as a comparison for Spain.

\paragraph{Three questions.} We keep three questions apart throughout.
\begin{enumerate}
  \item \emph{Association.} How large is the conditional association between GDP growth
    and the change in the unemployment rate in each country, and how much do countries
    differ once estimation uncertainty and common European shocks are accounted for?
  \item \emph{Prediction.} Does that association, or the hierarchical pooling of
    countries, improve forecasts of unemployment relative to simple benchmarks?
  \item \emph{Causation.} Can any part of the association be given a structural
    interpretation: the effect on unemployment of an identified shock that moves output?
\end{enumerate}
A large and stable association need not predict well; a useful predictor need not be
causal; and a causal effect of one kind of shock need not equal the reduced-form Okun
coefficient, which averages over all shocks of a period (section~\ref{sec:background}).

\paragraph{Approach.} Each stage was registered before its own estimation, sequentially,
on the same data: its protocol was committed and pushed before it was run, and its change
log records any later clarification with its reason. Later stages were designed knowing
the results of earlier ones (Appendix~\ref{app:prereg}). The data are frozen Eurostat and OECD snapshots
whose raw responses are hashed and re-derived by an audit. Every number in this article is
read by code from the committed result tables (Appendix~\ref{app:repro}). The work
proceeded in stages:
\begin{itemize}
  \item classical regressions and an independent review of them;
  \item a hierarchical Bayesian model with its checks and amendments;
  \item two predictive comparisons;
  \item an identification study with four strategies;
  \item a case study of Portugal, Spain, the Netherlands and Luxembourg.
\end{itemize}

\paragraph{Findings in brief.}
\begin{itemize}
  \item Outside 2020--2021 the two-lag sum is negative in \val{hi_point_negative} of
    \val{n_countries} countries and significantly so in \val{hi_negative} under
    small-sample bootstraps. How many individual sums survive an adjustment for testing
    \val{n_countries} countries depends on the error rate controlled.
  \item Its magnitude depends on the sample, the error model and the treatment of
    persistence, and country rankings are not established.
  \item Over all target quarters, no registered forecast comparison is decided. Without
    the pandemic targets, pooling, GDP lags and the common factor give small gains, but
    nothing reliably beats a forecast of no change.
  \item With the available public data, the instruments and restrictions we could defend
    are uninformative about a causal Okun coefficient.
  \item Of the three earlier claims, Spain's is supported, while the Dutch and Luxembourg
    claims are qualified. The commuter explanation for Luxembourg is not supported. The
    employment-protection correlations all have the opposite sign, but none is decided
    in the primary variant.
\end{itemize}

\paragraph{Outline.} Section~\ref{sec:background} sets out Okun's law and the accounting
identity behind it. Section~\ref{sec:data} describes the data and their measurement
issues. Sections~\ref{sec:classical} and~\ref{sec:bayes} present the classical and
Bayesian methods and estimates. Section~\ref{sec:heterogeneity} turns to uncertainty and
heterogeneity, section~\ref{sec:prediction} to prediction and
section~\ref{sec:identification} to identification. Section~\ref{sec:cases} contains the
country case studies and the comparison with the earlier claims, followed by the
discussion and conclusions. The appendices document reproduction and the pre-registration
record.
